%%%%%%%%%%%%%%%%%
% CV tailored for EDHEC Climate Institute - Senior Researcher Real Estate Climate Risk
% Positioning: Senior Researcher / Quantitative Climate Risk | Python, modelling, uncertainty, physical risk, reproducible pipelines & applied research
%%%%%%%%%%%%%%%%%

\documentclass[8pt,a4paper,ragged2e,withhyper]{altacv}

\geometry{left=1.25cm,right=1.25cm,top=.5cm,bottom=1.cm,columnsep=1.2cm}

\usepackage{paracol}
\usepackage{fontawesome5}
\usepackage{hyperref}

\iftutex
  \setmainfont{Roboto Slab}
  \setsansfont{Lato}
  \renewcommand{\familydefault}{\sfdefault}
\else
  \usepackage[rm]{roboto}
  \usepackage[defaultsans]{lato}
  \renewcommand{\familydefault}{\sfdefault}
\fi

\definecolor{SlateGrey}{HTML}{2E2E2E}
\definecolor{LightGrey}{HTML}{666666}
\definecolor{DarkPastelRed}{HTML}{8AC0D9}
\definecolor{PastelRed}{HTML}{00B0DA}
\definecolor{GoldenEarth}{HTML}{B4AD0C}

\colorlet{name}{black}
\colorlet{tagline}{PastelRed}
\colorlet{heading}{DarkPastelRed}
\colorlet{headingrule}{GoldenEarth}
\colorlet{subheading}{PastelRed}
\colorlet{accent}{PastelRed}
\colorlet{emphasis}{SlateGrey}
\colorlet{body}{LightGrey}

\definecolor{violet}{HTML}{5A2582}
\definecolor{rouge}{HTML}{F53B3B}
\definecolor{noir}{HTML}{00232C}
\definecolor{bleu}{HTML}{00B0DA}
\definecolor{rose}{HTML}{F831A0}
\definecolor{jaune}{HTML}{FFEC00}
\definecolor{vert}{HTML}{34AD6C}

\renewcommand{\namefont}{\Huge\rmfamily\bfseries}
\renewcommand{\personalinfofont}{\footnotesize}
\renewcommand{\cvsectionfont}{\LARGE\rmfamily\bfseries}
\renewcommand{\cvsubsectionfont}{\large\bfseries}
\renewcommand{\cvItemMarker}{{\small\textbullet}}
\renewcommand{\cvRatingMarker}{\faCircle}

\input{../../_shared/pubs-num.tex}

% auto_tex_manager layout tuning
\emergencystretch=3em
\hfuzz=60pt
\hbadness=9999
\sloppy

\begin{document}

\name{BOILEAU Guillaume}
\tagline{Senior Researcher / Quantitative Climate Risk | Python, modelling, uncertainty, physical risk \& applied research}

\photoL{2.8cm}{../../_shared/moi}

\personalinfo{%
  \email{guillaume.boileaupro@gmail.com}
  \phone{+33 6 59 94 85 84}
  \location{Alpes-Maritimes, FRANCE}
  \linkedin{guillaume-boileau-891b81265}
  \researchgate{Guillaume-Boileau-2}
  \orcid{0000-0002-3576-6968}
}

\makecvheader
\vspace{-0.3cm}

\AtBeginEnvironment{itemize}{\small}
\columnratio{0.64}

\begin{paracol}{2}

\cvsection{Experience}

\cvevent{\textbf{Software Engineer -- Data workflows, operational risk \& technical documentation}}{VEV Platform Services France}{2025 -- 2026}{Valbonne, France}

\textbf{Context:} Industrial software platform connected to electric vehicle charging infrastructure, with operational data flows, field constraints, incident analysis and digital service reliability.

\textbf{Objective:} Contribute to the reliability of operational digital services by analysing data flows, incidents, application behaviours and the interface between software, infrastructure and field operations.

\begin{itemize}
  \item \textbf{Operational data analysis:} Investigation of FTP data exchanges, interrupted flows, data inconsistencies and unexpected application behaviour.
  \item \textbf{Risk and incident analysis:} Diagnosis of operational incidents, root-cause analysis, formalisation of findings and support to corrective actions.
  \item \textbf{Infrastructure-related systems:} Work at the interface between software, field infrastructure, operational data and service constraints.
  \item \textbf{Technical documentation:} Production of clear technical notes, anomaly analyses and structured information for engineering teams.
  \item \textbf{Software contribution:} Development contributions in TypeScript, Angular and Node.js within an operational platform.
  \item \textbf{Collaborative tools:} Use of Linux, Git, Zenhub and technical documentation workflows.
\end{itemize}

\divider

\cvevent{\textbf{Senior R\&D Engineer -- Quantitative modelling, uncertainty \& reproducible pipelines}}{CNES}{2023 -- 2025}{Nice, France}

\textbf{Context:} Applied R\&D and software development for the LISA ESA/NASA space mission ground segment, involving mission data chains L0--L3, simulation, model integration, uncertainty, validation and reproducible scientific workflows.

\textbf{Objective:} Design robust Python workflows to integrate, compare, validate and document complex models and heterogeneous data, with emphasis on uncertainty, traceability, reproducibility and methodological rigour.

\begin{itemize}
  \item \textbf{Quantitative research workflows:} Development of CATpyLISA, a Python platform for model integration, comparison, validation and technical review.
  \textcolor{DarkPastelRed}{\href{https://gitlab.oca.eu/gboileau/lisa_l3_tools}{\bf GitLab: CATpyLISA}}
  \item \textbf{Methodology development:} Structuring of assumptions, input datasets, configurations, intermediate results, model outputs and validation metrics.
  \item \textbf{Uncertainty and sensitivity:} Analysis of weak signals, noise contributions, parameter effects, limitations, model behaviour and robustness.
  \item \textbf{Model validation:} Definition of acceptance criteria, comparison rules, consistency checks, validation procedures and gap analyses.
  \item \textbf{Reproducible pipelines:} Design of documented workflows enabling traceable analyses, reviewability and collaborative use.
  \item \textbf{Applied research outputs:} Preparation of technical notes, review material, methodological explanations and decision-support summaries.
  \item \textbf{International collaboration:} Interface with scientific, software and institutional contributors in the LISA international context.
\end{itemize}

\divider

\cvevent{\textbf{Data Scientist -- Physical risk proxies, sensor data \& statistical modelling}}{Universiteit Antwerpen, Belgium}{2021 -- 2023}{Antwerp, Belgium}

\textbf{Context:} Applied data science for precision instruments and next-generation detectors, combining sensor data, environmental noise, statistical modelling, machine learning and performance validation.

\textbf{Objective:} Develop robust workflows to analyse sensor datasets, characterise environmental contributions, model uncertainties, compare configurations and produce decision-oriented recommendations.

\begin{itemize}
  \item \textbf{Sensor and environmental data:} Analysis of environmental measurements, spatial correlations, propagation effects and their impact on system performance.
  \item \textbf{Physical risk reasoning:} Assessment of how external physical phenomena propagate into system-level performance degradation and uncertainty.
  \item \textbf{Statistical modelling:} Development of MCMC pipelines for different configurations and noise assumptions.
  \item \textbf{Machine learning:} Development of approaches for non-stationary noise reduction using witness channels, machine learning and deep learning.
  \item \textbf{Performance metrics:} Figures of merit, sensitivity analyses, consistency checks and comparison of configurations.
  \item \textbf{Research communication:} Production of traceable results, technical syntheses and recommendations for international stakeholders.
\end{itemize}

\divider

\cvevent{\textbf{Junior R\&D Engineer -- Scientific modelling, simulation \& publications}}{ARTEMIS, Observatoire de la Côte d’Azur}{2018 -- 2021}{Nice, France}

\textbf{Context:} Applied research on weak-signal analysis, numerical simulation, signal processing and preparation of space mission science cases.

\textbf{Objective:} Develop modelling and simulation methods to produce robust, documented and publishable scientific results.

\begin{itemize}
  \item \textbf{Scientific modelling:} Development of numerical simulation workflows for complex signal environments and large datasets.
  \item \textbf{Signal processing:} Methods to analyse and separate weak, overlapping signals.
  \item \textbf{Validation:} Consistency checks, result comparison, sensitivity studies and methodological limitation analysis.
  \item \textbf{Academic writing:} Contribution to scientific publications, reports, methodological explanations and presentations.
\end{itemize}

\switchcolumn

\cvsection{Profile}

\textbf{PhD-level quantitative researcher with strong experience in applied modelling, uncertainty analysis, Python pipelines, reproducible research, physical-risk reasoning, model validation and scientific publication.}

Target position: contribute to EDHEC Climate Institute’s real estate climate risk methodology by bringing rigorous quantitative modelling, transparent assumptions, documented pipelines, uncertainty analysis and a strong capacity to bridge research outputs with operational decision-making.

\begin{itemize}
  \item Strong research foundation: PhD in physics, peer-reviewed scientific output, applied R\&D and international collaboration.
  \item Robust quantitative toolkit: Python, statistical modelling, simulation, model validation, uncertainty, sensitivity analysis and reproducible pipelines.
  \item Relevant physical-risk mindset: experience modelling how external physical phenomena, environmental noise and system vulnerabilities affect performance.
  \item Strong writing ability for scientific, technical and practitioner-oriented audiences.
  \item Targeted ramp-up on real estate finance, climate finance, CRREM, EPC datasets, LTV, PD/LGD, SFDR, CSRD and EBA guidelines.
\end{itemize}

\cvsection{Languages}

\cvskill{English}{4}
\divider
\cvskill{Spanish}{2}
\divider

\medskip

\cvsection{Education}

\cvevent{PhD in Physics}{Université Côte d’Azur}{2018 -- 2021}{Nice, France}

\divider

\cvevent{Selective Graduate Program in Fundamental Physics (Magistère)}{Université Paris-Saclay}{2016 -- 2018}{Orsay, France}

\divider

\cvevent{MSc in Astronomy and Astrophysics}{Observatoire de Paris / Université Paris-Saclay}{2017 -- 2018}{Paris, France}

\divider

\cvevent{BSc in Fundamental Physics}{Université Paris-Saclay}{2015 -- 2016}{Orsay, France}

\cvsection{Professional Training}

\cvevent{Machine Learning in Python with scikit-learn}{FUN MOOC -- Inria}{2026}{France Université Numérique}

\small{
Predictive modelling, supervised learning, model evaluation, cross-validation, metrics, pipelines and methodological/software fundamentals of machine learning in Python.
}

\divider

\cvevent{Reproducible Research: Methodological Principles for Transparent Science}{FUN MOOC -- Inria}{2025}{France Université Numérique}

\small{
Reproducible research, GitLab, notebooks/Jupyter, data management, computational documentation, software environments and reproducibility methods.
}

\divider

\cvevent{Continuous Professional Training -- Data, AI, Software and Systems}{OpenClassrooms / Coursera / FUN MOOC}{2020 -- 2026}{}

\small{
Python, SQL, Machine Learning, AI, Linux, JavaScript, TypeScript, Angular, Node.js, Django, REST APIs, C/C++, data analysis and software engineering fundamentals.
}

\end{paracol}

\cvsection{Projects}

\cvevent{CATpyLISA -- Reproducible modelling and validation platform}{Python / Scientific workflows / Model comparison}{2023 -- 2025}{}

Python platform for integrating models, structuring data, comparing results, validating outputs and supporting collaborative technical reviews in a complex scientific environment.

\textbf{Role:} Python development, model validation, data structuring, reproducible workflows, sensitivity checks, gap analysis, documentation and coordination with multiple contributors.

\textbf{Transfer to climate risk:} methodology development, transparent assumptions, reproducible pipelines, model validation, uncertainty analysis and practitioner-oriented documentation.

\divider

\cvevent{Environmental data, risk propagation and performance modelling}{Sensor data / Physical risk / Statistical modelling}{2021 -- 2023}{}

Development of numerical and statistical workflows to analyse environmental sensor data, characterise physical disturbances, compare configurations and assess performance impacts.

\textbf{Role:} sensor data analysis, statistical modelling, MCMC pipelines, non-stationary noise analysis, performance metrics, validation and reporting.

\textbf{Transfer to real estate climate risk:} reasoning from physical shocks to system-level vulnerability, uncertainty, scenario comparison and decision-support metrics.

\divider

\cvevent{Scientific publications and applied research outputs}{Peer-reviewed research / Modelling / Simulation}{2018 -- 2025}{}

Scientific work involving modelling, simulation, uncertainty, validation, weak-signal analysis and peer-reviewed publication in international research environments.

\textbf{Role:} modelling, numerical analysis, scientific writing, result validation, interpretation, communication and collaborative review.

\textbf{Transfer to EDHEC Climate Institute:} academic rigour, ability to publish, build intellectual presence and translate quantitative frameworks into usable outputs.

\cvsection{Core Skills}

\vspace{-.3cm}

\cvsubsection{Quantitative research \& modelling}

{\small
\cvtag{Quantitative modelling}
\cvtag{Applied research}
\cvtag{Methodology development}
\cvtag{Uncertainty analysis}
\cvtag{Sensitivity analysis}
\cvtag{Scenario analysis}
\cvtag{Statistical modelling}
\cvtag{MCMC}
\cvtag{Model validation}
\cvtag{Assumption testing}
\cvtag{Boundary conditions}
\cvtag{Scientific publication}
}

\divider\smallskip
\vspace{-.3cm}

\cvsubsection{Climate risk \& real estate finance}

{\small
\cvtag{Climate risk -- ramp-up}
\cvtag{Physical risk}
\cvtag{Transition risk -- ramp-up}
\cvtag{Real estate finance -- ramp-up}
\cvtag{DCF -- ramp-up}
\cvtag{Income capitalisation -- ramp-up}
\cvtag{LTV -- ramp-up}
\cvtag{PD/LGD -- ramp-up}
\cvtag{Stranded assets -- awareness}
\cvtag{Insurance availability -- awareness}
\cvtag{Retrofit CapEx -- awareness}
\cvtag{Regulatory costs -- awareness}
}

\divider\smallskip
\vspace{-.3cm}

\cvsubsection{Built environment data \& frameworks}

{\small
\cvtag{EPC datasets -- ramp-up}
\cvtag{Energy labels -- ramp-up}
\cvtag{Building stock models -- ramp-up}
\cvtag{Building archetypes -- ramp-up}
\cvtag{CRREM -- ramp-up}
\cvtag{LCA frameworks -- ramp-up}
\cvtag{Data coverage}
\cvtag{Data quality}
\cvtag{Calibration strategies}
\cvtag{Substitution strategies}
\cvtag{Portfolio datasets -- ramp-up}
}

\divider\smallskip
\vspace{-.3cm}

\cvsubsection{Python, data \& reproducibility}

{\small
\cvtag{Python}
\cvtag{pandas}
\cvtag{NumPy}
\cvtag{SciPy}
\cvtag{Matplotlib}
\cvtag{Jupyter}
\cvtag{SQL}
\cvtag{Git / GitLab}
\cvtag{Linux}
\cvtag{Reproducible pipelines}
\cvtag{Documentation}
\cvtag{Data analysis}
\cvtag{Data validation}
}

\divider\smallskip
\vspace{-.3cm}

\cvsubsection{Regulation, finance stakeholders \& communication}

{\small
\cvtag{SFDR -- awareness}
\cvtag{CSRD -- awareness}
\cvtag{EBA guidelines -- awareness}
\cvtag{Financial institutions -- ramp-up}
\cvtag{Regulators -- awareness}
\cvtag{Client-facing notes}
\cvtag{Academic writing}
\cvtag{Policy forums -- ramp-up}
\cvtag{Industry workshops}
\cvtag{Cross-disciplinary collaboration}
\cvtag{English}
}

\divider\smallskip
\vspace{-.3cm}

\cvsubsection{Systems, risk \& decision support}

{\small
\cvtag{Risk transmission}
\cvtag{Vulnerability modelling}
\cvtag{Portfolio-level reasoning}
\cvtag{Performance metrics}
\cvtag{Decision-support tools}
\cvtag{Complex systems}
\cvtag{Physical processes}
\cvtag{Validation}
\cvtag{Robustness}
\cvtag{Technical synthesis}
\cvtag{Stakeholder communication}
}

\end{document}
